% LSPV-4 REPORT::PCLOSE1-R1::2026-10-01T13:43:00-07:00
\documentclass[11pt]{article}
\usepackage[margin=0.85in]{geometry}
\usepackage{amsmath,amssymb}
\begin{document}
\section*{LSPV-4 hostile audit: P-closure of polynomial-depth CCSPC}
\textbf{STATUS: OPEN.} P versus NP remains open. This report proves no unrestricted SAT lower bound.
\paragraph{CORPUS REVIEWED.} I surveyed root, tasks, current/historical MAIN, LSPV-1--5 reports bearing on SU/vMS, reconstruction, acquisition and causal branches, governance, and HEC/CRS--PPC/CCSPC artifacts. I read AGENTS.md, CODEX.md, pnp-research skill, root main.tex, TASK-004--007, latest synthesis surfaced by search, LSPV-2 HADSQ1, LSPV-3 DEPSWAP1, LSPV-4 ACQ1/REREAD1, CCSPC1, CRS--PPC best theorems, and HEC proved core. Predecessor DEPSWAP1 exists and is substantive. Repository search located prior pointer-jumping, reread, group-compression, NEXP-guard and barrier claims.
\paragraph{Two obstacles and literature preflight.} First is the static arbitrary-algorithm bridge. Second is closure under the counterfactual: if every CCSPC stage is polynomial once SAT is polynomial, polynomially many adaptive stages remain polynomial. Mao--Yang--Zhang prove an \(Omega(n/k+k)\) restricted-round communication lower bound for pointer chasing; Fischer--Oshman--Rosen--Roth prove communication lower bounds even with unlimited interaction. These do not transfer to centralized static SAT. Aaronson--Wigderson's algebrization barrier still applies to oracle-stable proof lines.
\paragraph{Exact P-closure theorem.} Fix a staticized CCSPC instance with base input \(x\) and complete challenge string \(G\). Suppose stage count \(z(n)\) is polynomial. Suppose each codebook constructor \(C_i\), audit selector \(A_i\), audit evaluator \(V_i\), and pointer update \(F_i\) is deterministic polynomial time in ordinary explicit input/transcript length, and canonical SAT response uses polynomially many adaptive SAT decisions plus polynomial deterministic work. Under \(P=NP\), the complete staticized trajectory and final output are in deterministic polynomial time.
\paragraph{Proof.} Under \(P=NP\), fix a polynomial-time exact SAT decider. A canonical satisfying assignment such as lexicographically maximal SAT is recovered by prefix self-reduction with at most one SAT decision per variable; UNSAT is detected initially. Thus each canonical response is polynomial. Compute each remaining declared stage primitive and read the already-static challenge block from \(G\). Repeating polynomially many polynomial-time stages is polynomial. \(square\)
\paragraph{Exact circuit-composition corollary.} If deterministic maps \(F_i:{0,1}^r\to{0,1}^r\) have circuit sizes \(s_i\), the endpoint \(F_z\circ\cdots\circ F_1\) has a circuit of size at most \(\sum_i s_i+O(zr)\) by sequential wiring. Polynomially many polynomial-size stage maps therefore have a polynomial-size static endpoint representation. This does not imply logarithmic evaluation time, but causal depth alone cannot force superpolynomial representation size.
\paragraph{Hostile consequence.} Canonical SAT dependence, delayed audits, adaptive pointers, and exponentially many latent paths cannot alone yield the contradiction. If all staticized stage primitives remain polynomial under \(P=NP\), the cascade is polynomial. A successful route must force a resource not closed under polynomial composition, use superpolynomial realized depth without silently changing to succinct/NEXP input, or prove a model-preserving compiler to a separately established superpolynomial resource.
\paragraph{Pointer-jumping reconciliation.} REREAD1 already proves causal pointer dependence does not force persistent static storage; the corpus also records group compression and pointer jumping. The new point is stronger P-closure. Even when stage maps differ and binary lifting gives no \(O(\log z)\) shortcut, ordinary sequential composition remains polynomial for polynomial \(z\). Failure of a fast pointer-jumping shortcut does not approach a superpolynomial SAT lower bound.
\paragraph{SU/vMS, ACQ1, PQ.} HADSQ1 fixed-Hadamard rank formulas and ACQ1 conditional legal-state separation survive. Exact SU09/vMS label-bit invocation and global stacked rank remain OPEN. ACQ1 is compatible with \(P=NP\): once legality and \(B_vq\neq0\) are polynomial predicates, the assumed SAT decider can perform NP search. Schur products, multiplication kernels, puncturing, fixed-code Hilbert functions and syzygies are EXACT transfers after the labelled binary realization is fixed. Succinct shortening/syzygy separation remains ANALOGY/CANDIDATE. Cryptographic hardness/random-oracle assumptions remain inadmissible.
\paragraph{Oracle and wrapper audit.} The P-closure theorem relativizes: semantic SAT oracles compose stagewise, and algebraic stages remain visible to algebraic oracles. It gives no nonalgebrizing hinge. DEAD/COPY/DELAY inflate traces without changing endpoints; INTERP/RECOMP can replace implementations. Horn, 2-SAT, bounded-treewidth and Gaussian/Tseitin controls show long dependencies can coexist with efficient global algorithms.
\paragraph{DUPLICATION / CONTRADICTION AUDIT.} CCSPC1 already says polynomially many adaptive NP-oracle computations collapse under \(P=NP\) and warns against NEXP conclusions from latent branching. New here is an explicit closure theorem for the whole staticized cascade, including canonical search, delayed audit and pointer update, plus the circuit-composition corollary. No contradiction with restricted HEC/CRS--PPC or communication pointer chasing: resources differ.
\paragraph{STALE OR OVERSTATED MAIN CLAIMS.} Root main.tex correctly leaves static-to-CRS--PPC and the nonrelativizing/nonalgebrizing hinge OPEN, but its causal discussion predates REREAD1 and omits full P-closure. Any reading that causal sequentiality or canonicalization alone forces superpolynomial static work is stale.
\paragraph{ADOPTED / REJECTED / RETIRED.} Adopt exact HEC/CRS freshness and restricted no-reread bounds, HADSQ1 fixed-encoder algebra, ACQ1 charged acquisition, and explicit-input normalization. Reject communication pointer-chasing transfer without a compiler. Retire causal depth alone, latent-path count, canonicalization alone, and polynomial-depth adaptive SAT composition as sources of superpolynomial contradiction.
\paragraph{COMPONENTS CONSISTENT WITH P=NP.} Static reread, polynomially many adaptive SAT decisions, canonical SAT self-reduction, stagewise nondeterminism, memoization, group compression, polynomial circuit composition and ACQ1 NP-search acquisition are compatible with \(P=NP\). This is not evidence for \(P=NP\).
\paragraph{Strict status.} \textbf{PROVED/EXACT:} P-closure theorem under stated hypotheses; polynomial circuit composition; REREAD1 bypass; HADSQ1 fixed-Hadamard formulas. \textbf{KNOWN THEOREM:} communication pointer-chasing bounds. \textbf{CONDITIONAL:} ACQ1 after exact evaluator transcription/compilation. \textbf{CANDIDATE:} representation-sensitive sequential work and succinct algebraic separation. \textbf{FALSIFIED/RETIRED:} causal depth alone as superpolynomial static lower bound; latent branching as runtime; generic communication-to-SAT transfer. \textbf{OPEN:} exact SU09/vMS invocation, stacked Schur rank/common kernel, ACQ1 exponent, sparse-Q normalization, reconstruction ledger, arbitrary-algorithm transfer, CRS--PPC/CCSPC compiler, nonrelativizing/nonalgebrizing hinge, final contradiction. P versus NP remains OPEN.
\paragraph{Weakest repair and next falsification target.} Formulate a CCSPC resource not automatically closed under polynomial sequential composition and prove a model-preserving compiler forcing it for every static SAT decider. Next instantiate CCSPC-E and replace every canonical SAT stage by the assumed polynomial decider plus self-reduction. If the end-to-end simulation is polynomial, that compiler does not create a contradiction.
\paragraph{Primary literature checked.} Mao, Yang, Zhang, Gadgetless Lifting Beats Round Elimination, arXiv:2411.10996. Fischer, Oshman, Rosen, Roth, Pointer Chasing with Unlimited Interaction, arXiv:2508.19158. Aaronson and Wigderson, Algebrization: A New Barrier in Complexity Theory, ACM TOCT 1(1), 2009.
\end{document}
% END_OF_VERIFIED_REPORT::LSPV-4::PCLOSE1-R1